\chapter{Kommentierte Quellen}

Achtundzwanzig Werke, eingeordnet: Zu jedem steht, wofür man es aufschlägt und was es leistet, das andere nicht leisten.

\lit{1}{Atiyah, M., Macdonald, I.: \textit{Introduction to Commutative Algebra}. Addison-Wesley 1969.}{Die kommutative Algebra hinter Teil~I und~III auf $128$ Seiten: Ideale, Radikale, Primärzerlegung, Noether-Normalisierung, Dimensionstheorie. \textbf{Die Fundstelle für die Auslassungen~1, 2 und~6}; knapp, mit Übungen, die die halbe Theorie tragen. Für Leser mit Band 11.}

\lit{2}{Bellman, R., Åström, K.\ J.: \textit{On Structural Identifiability}. Math.\ Biosci.\ 7, 1970.}{Der Aufsatz, der das Problem des Leitfalls stellt und benennt: Welche Parameter eines Kompartimentmodells bestimmen die Daten? \textbf{Sieben Seiten, aus denen ein Gebiet wurde}; die Methode ist die Ein-Ausgabe-Gleichung aus Kapitel~6, ohne Gröbner-Basen.}

\lit{3}{Brieskorn, E., Knörrer, H.: \textit{Ebene algebraische Kurven}. Birkhäuser 1981.}{\textbf{Auf Deutsch, siebenhundert Seiten über ebene Kurven}: Geschichte von den Griechen bis Newton, projektive Ebene, Bézout, Singularitäten, Puiseux, Geschlecht --- mit mehr Bildern als jedes andere Buch des Fachs. Die Kapitel~9, 11 und~12 in größter Breite; die Fundstelle für Auslassung~17 und für jede Kurve, die hier nur erwähnt wurde.}

\lit{4}{Bosch, S.: \textit{Algebraische Geometrie und kommutative Algebra}. 2.\ Aufl., Springer Spektrum 2020.}{\textbf{Auf Deutsch}: der Weg von den Idealen zu den Schemata, mit vollständigen Beweisen --- Nullstellensatz, Dimension, Glattheit, dann Garben und Schemata. Die Fortsetzung dieses Bandes in die Sprache Grothendiecks, für Leser, die Kapitel~15 ernst nehmen.}

\lit{5}{Cox, D., Little, J., O'Shea, D.: \textit{Ideals, Varieties, and Algorithms}. 4.\ Aufl., Springer 2015.}{\textbf{Das Lehrbuch, dem dieser Band in Teil~I bis III am meisten verdankt}: Gröbner-Basen vor der Theorie, Elimination, Dimension über Hilbert-Funktionen, Robotik als Anwendung. Die Fundstelle für die Auslassungen~3, 5, 6, 10 und~11; mit Programmieraufgaben, die das Programm dieses Bandes ergänzen.}

\lit{6}{Eisenbud, D.: \textit{Commutative Algebra with a View Toward Algebraic Geometry}. Springer 1995.}{Achthundert Seiten kommutative Algebra, jede mit einem geometrischen Grund: Primärzerlegung, Dimension, Gröbner-Basen (Kapitel~15), Syzygien, Kohen--Macaulay. \textbf{Die Algebra zu Hartshorne}, und das Buch, das man aufschlägt, wenn Atiyah--Macdonald zu knapp ist.}

\lit{7}{Fulton, W.: \textit{Algebraic Curves: An Introduction to Algebraic Geometry}. Benjamin 1969; frei verfügbar 2008.}{Ebene Kurven mit allen Beweisen auf $120$ Seiten: lokale Ringe, Schnittvielfachheit, Bézout, Cayley--Bacharach, Riemann--Roch. \textbf{Die Fundstelle für die Auslassungen~12, 14, 17 und~18}; der Autor hat das Buch 2008 frei ins Netz gestellt. Für Teil~IV die erste Adresse.}

\lit{8}{Gelfand, I., Kapranov, M., Zelevinsky, A.: \textit{Discriminants, Resultants, and Multidimensional Determinants}. Birkhäuser 1994.}{Die Rückkehr der Eliminationstheorie: Resultanten und Diskriminanten für beliebig viele Gleichungen, mit Newton-Polytopen und Hypergeometrie. \textbf{Schwer, grundlegend, der Ursprung der dünnen Resultanten}; Kapitel~6 und~9 in ihrer modernen Fortsetzung.}

\lit{9}{Greuel, G.-M., Pfister, G.: \textit{A Singular Introduction to Commutative Algebra}. 2.\ Aufl., Springer 2008.}{Kommutative Algebra mit dem Rechner, geschrieben von den Autoren des Systems Singular in Kaiserslautern: jeder Satz mit einem Beispiel, das man eintippen kann. \textbf{Das Buch für Leser, die nach Kapitel~5 wirklich rechnen wollen} --- Standardbasen, Primärzerlegung, Hilbert-Funktionen, Normalisierung, alles implementiert.}

\lit{10}{Hartshorne, R.: \textit{Algebraic Geometry}. Springer 1977.}{Das Standardwerk: Kapitel~I die klassischen Varietäten (die Fundstelle für die Auslassungen~4, 7, 8, 9 und~16), Kapitel~II und~III Schemata und Kohomologie, Kapitel~IV Kurven, Kapitel~V Flächen. \textbf{Schwer, aber nach diesem Band ist Kapitel~I lesbar}; die Übungen enthalten die halbe Geometrie.}

\lit{11}{Hilbert, D.: \textit{Ueber die Theorie der algebraischen Formen}. Math.\ Ann.\ 36, 1890.}{Basissatz, Syzygiensatz und Hilbert-Polynom in einer Arbeit --- die Grundlage von Kapitel~2 und~7 im Original. \textbf{Die Arbeit, die Gordan Theologie nannte}; der Beweis ist anders als der hier gegebene und lesenswert als Dokument, wie Mathematik ihre Methoden wechselt.}

\lit{12}{Hulek, K.: \textit{Elementare Algebraische Geometrie}. 2.\ Aufl., Vieweg+Teubner 2012.}{\textbf{Auf Deutsch, für das vierte Semester}: affine und projektive Varietäten, Dimension, Glattheit, Kurven, Kubikflächen --- die Teile~I, III und~IV dieses Bandes in strenger Form, mit den Beweisen, die hier Auslassungen sind. Die deutschsprachige Entsprechung zu Reid und Perrin.}

\lit{13}{Kunz, E.: \textit{Einführung in die algebraische Geometrie}. Vieweg 1997.}{\textbf{Auf Deutsch, aus Regensburg}: der algebraische Zugang --- Nullstellensatz, Dimension, reguläre Punkte, Divisoren auf Kurven --- mit vollständigen Beweisen und einer Klarheit, die ihn zum besten deutschsprachigen Text für die Auslassungen~1, 6 und~9 macht. Kunz' \enquote{Einführung in die kommutative Algebra} (1980) ist der Begleitband.}

\lit{14}{Meshkat, N., Sullivant, S., Eisenberg, M.: \textit{Identifiability Results for Several Classes of Linear Compartment Models}. Bull.\ Math.\ Biol.\ 77, 2015.}{Der Leitfall in voller Allgemeinheit: lineare Kompartimentmodelle als Varietäten, Identifizierbarkeit über die Ein-Ausgabe-Gleichung und ihre Koeffizienten, Sätze über Ketten, Zyklen und Bäume. \textbf{Die Quelle für Kapitel~13.1} und der Beweis, dass die Methode dieses Bandes die der Praxis ist.}

\lit{15}{Mumford, D.: \textit{The Red Book of Varieties and Schemes}. Springer 1988 (Vorlesung 1967).}{Die Motivation für Schemata, von dem, der sie am besten erklären konnte: warum $V(x^{2})$ etwas anderes ist als $V(x)$, warum man Primideale als Punkte braucht, mit Mumfords berühmten Zeichnungen. \textbf{Das Buch für die Frage, die Kapitel~15 offenlässt}; vor Hartshorne zu lesen.}

\lit{16}{Pachter, L., Sturmfels, B. (Hg.): \textit{Algebraic Statistics for Computational Biology}. Cambridge 2005.}{Statistische Modelle als Varietäten, phylogenetische Invarianten, Markov-Basen, tropische Geometrie --- das Buch, das die algebraische Statistik der Biologie bekannt machte. \textbf{Kapitel~13 in der Fassung für Genetiker}; mit Beispielen, die man mit dem Programm dieses Bandes nachrechnen kann.}

\lit{17}{Perrin, D.: \textit{Algebraic Geometry: An Introduction}. Springer 2008.}{Die klassische Theorie mit Blick auf Schemata, in einer Vorlesung aus Paris: Varietäten, Garben, Dimension, Tangentialräume, Bézout, Kurven, Riemann--Roch. \textbf{Zwischen Reid und Hartshorne}; die Fundstelle für Auslassung~18 in einer Fassung, die ohne Kohomologie auskommt.}

\lit{18}{Raghavan, M., Roth, B.: \textit{Inverse Kinematics of the General 6R Manipulator and Related Linkages}. J.\ Mech.\ Design 115, 1993.}{Die Gleichung vom Grad $16$ für den allgemeinen Sechs-Gelenk-Arm, hergeleitet durch Elimination --- Kapitel~14 im Original. \textbf{Der Aufsatz, der die Robotik zur Anwenderin der Eliminationstheorie machte}; mit der Geschichte der Zahl $16$ in der Einleitung.}

\lit{19}{Reid, M.: \textit{Undergraduate Algebraic Geometry}. Cambridge 1988.}{Hundertdreißig Seiten, die zeigen, dass algebraische Geometrie im dritten Studienjahr lehrbar ist: Kegelschnitte, Kubiken mit Gruppengesetz, Nullstellensatz, Dimension, Singularitäten, die $27$ Geraden. \textbf{Der nächste Verwandte dieses Bandes im Ton}, mit einem Schlusskapitel über die Geschichte des Fachs, das man gelesen haben sollte.}

\lit{20}{Schenck, H.: \textit{Computational Algebraic Geometry}. Cambridge 2003.}{Gröbner-Basen, Hilbert-Funktionen, Syzygien und Kohomologie mit Macaulay2 --- der rechnende Weg von Kapitel~5 in die Kohomologie, mit Code auf jeder zweiten Seite. \textbf{Für Leser, die Cox--Little--O'Shea gelesen haben und wissen wollen, was dahinter kommt.}}

\lit{21}{Sendra, J.\ R., Winkler, F., Pérez-Díaz, S.: \textit{Rational Algebraic Curves: A Computer Algebra Approach}. Springer 2008.}{Kapitel~11 als Algorithmus: Geschlecht berechnen, Parametrisierung finden, über $\mathbb Q$ oder gar nicht, mit Beweisen und Implementierung. \textbf{Die Fundstelle für Auslassung~16 und für alles, was \enquote{parametrize} in einem Computeralgebrasystem tut.}}

\lit{22}{Shafarevich, I.: \textit{Basic Algebraic Geometry 1: Varieties in Projective Space}. 3.\ Aufl., Springer 2013.}{Die klassische Theorie vollständig: Dimension, Glattheit, Divisoren, Schnitttheorie, mit der Anschauung eines Zahlentheoretikers. \textbf{Die Fundstelle für die Auslassungen~7, 9 und~11}; Band~2 behandelt Schemata und komplexe Mannigfaltigkeiten. Länger als Hartshorne Kapitel~I und leichter.}

\lit{23}{Silverman, J.: \textit{The Arithmetic of Elliptic Curves}. 2.\ Aufl., Springer 2009.}{Das Standardwerk zu Kapitel~10: Gruppengesetz, Isogenien, Kurven über endlichen Körpern (Hasse, Schoof), über lokalen Körpern, Mordell--Weil mit Beweis, Höhen. \textbf{Die Fundstelle für Auslassung~15}; schwer, vollständig, unersetzt.}

\lit{24}{Silverman, J., Tate, J.: \textit{Rational Points on Elliptic Curves}. 2.\ Aufl., Springer 2015.}{Dasselbe Thema für Studierende: Gruppengesetz, Nagell--Lutz mit Beweis, Mordell für Kurven mit rationalem Zwei-Torsionspunkt \textbf{vollständig bewiesen auf dreißig Seiten}, Kurven über endlichen Körpern, Faktorisierung mit elliptischen Kurven. Das Buch, mit dem man Kapitel~10 vertieft, bevor man Silverman aufschlägt.}

\lit{25}{Sturmfels, B.: \textit{Solving Systems of Polynomial Equations}. AMS 2002.}{Zehn Vorlesungen über das Lösen: Gröbner-Basen und Eigenwerte, Resultanten, Bernstein--Kushnirenko--Khovanskii, numerische Homotopie, reelle Lösungen, Anwendungen in Statistik und Robotik. \textbf{Die Fundstelle für Auslassung~13} und das beste Buch über die Frage, wie viele Lösungen ein System hat.}

\lit{26}{Sullivant, S.: \textit{Algebraic Statistics}. AMS 2018.}{Das Lehrbuch der algebraischen Statistik: Modelle als Varietäten, Identifizierbarkeit, Mischungen und Sekantenvarietäten, Markov-Basen, Maximum-Likelihood-Grad, phylogenetische Modelle. \textbf{Kapitel~13 auf fünfhundert Seiten}; die Fundstelle für Auslassung~20, von einem der Autoren von \zit{14}.}

\lit{27}{Washington, L.: \textit{Elliptic Curves: Number Theory and Cryptography}. 2.\ Aufl., Chapman \& Hall 2008.}{Elliptische Kurven von der Addition bis zu ECDSA, Paarungen und Angriffen (MOV, anomale Kurven), mit Zahlenbeispielen, die man nachrechnen kann. \textbf{Die Quelle für Kapitel~14.3 bis 14.5}; für Leser mit Band 1 und Band 20.}

\lit{28}{Zariski, O., Samuel, P.: \textit{Commutative Algebra}, 2 Bände. Van Nostrand 1958/60.}{Die kommutative Algebra, wie Zariski sie für die Geometrie brauchte: Körpererweiterungen, Lüroth, Noether-Normalisierung, Dimension, Bewertungen. \textbf{Das Werk, das die italienische Geometrie auf algebraischen Boden stellte}; als Nachschlagewerk bis heute unübertroffen, als Lektüre nur für Geduldige.}
